% myhill-nerode.tex — Myhill-Nerode congruence: futures equivalence partitions states
% Requires opoch.sty (loaded by main document)

\begin{figure}[ht]
\centering
\begin{tikzpicture}[
  state/.style={circle, draw=opochgray!80, fill=white, minimum size=9mm,
    font=\footnotesize, inner sep=1pt, line width=0.6pt},
  future arrow/.style={-{Stealth[length=2mm]}, thin, #1},
  outcome/.style={rounded rectangle, draw=#1, fill=#1!15, font=\scriptsize\bfseries,
    minimum width=10mm, minimum height=6mm, inner sep=2pt},
  class label/.style={font=\scriptsize\bfseries, #1},
]

  % === Equivalence Class 1 (blue): s1, s2, s3 ===
  \begin{scope}[on background layer]
    \draw[rounded corners=10pt, draw=opochblue!50, fill=opochblue!8,
      dashed, line width=0.8pt]
      (-0.8, -0.5) rectangle (4.8, 2.5);
  \end{scope}
  \node[class label=opochblue!80!black] at (2.0, 2.2) {$[s]_1$};

  \node[state] (s1) at (0, 1.0) {$s_1$};
  \node[state] (s2) at (2, 1.0) {$s_2$};
  \node[state] (s3) at (4, 1.0) {$s_3$};

  % === Equivalence Class 2 (green): s4, s5 ===
  \begin{scope}[on background layer]
    \draw[rounded corners=10pt, draw=opochgreen!50, fill=opochgreen!8,
      dashed, line width=0.8pt]
      (-0.8, -3.5) rectangle (2.8, -1.0);
  \end{scope}
  \node[class label=opochgreen!80!black] at (1.0, -1.3) {$[s]_2$};

  \node[state] (s4) at (0, -2.2) {$s_4$};
  \node[state] (s5) at (2, -2.2) {$s_5$};

  % === Equivalence Class 3 (orange): s6, s7, s8 ===
  \begin{scope}[on background layer]
    \draw[rounded corners=10pt, draw=opochorange!50, fill=opochorange!8,
      dashed, line width=0.8pt]
      (3.2, -3.5) rectangle (8.8, -1.0);
  \end{scope}
  \node[class label=opochorange!80!black] at (6.0, -1.3) {$[s]_3$};

  \node[state] (s6) at (4, -2.2) {$s_6$};
  \node[state] (s7) at (6, -2.2) {$s_7$};
  \node[state] (s8) at (8, -2.2) {$s_8$};

  % === Terminal outcomes (right side) ===
  \node[outcome=opochgreen] (unique) at (10.5, 1.0) {$\UNIQUE$};
  \node[outcome=opochorange] (omega)  at (10.5, -2.2) {$\OMEGA$};

  % === Futures arrows from Class 1 states -> UNIQUE ===
  \draw[future arrow=opochblue!70] (s1) -- ++(1.8, 0)
    node[font=\tiny, above, pos=0.4, opochblue!70!black] {$\sigma_1$};
  \draw[future arrow=opochblue!70] (s2) -- ++(1.8, 0)
    node[font=\tiny, above, pos=0.4, opochblue!70!black] {$\sigma_1$};
  \draw[future arrow=opochblue!70] (s3) -- ++(1.8, 0)
    node[font=\tiny, above, pos=0.4, opochblue!70!black] {$\sigma_1$};

  % Merge arrows toward UNIQUE
  \draw[future arrow=opochblue!70] (s1.east) ++(1.8, 0) -- (unique.west);
  \draw[future arrow=opochblue!70] (s2.east) ++(1.8, 0) -- (unique.west);
  \draw[future arrow=opochblue!70] (s3.east) ++(1.8, 0) -- (unique.west);

  % === Futures arrows from Class 2 states -> UNIQUE ===
  \draw[future arrow=opochgreen!70] (s4) -- ++(2.0, 0)
    node[font=\tiny, above, pos=0.4, opochgreen!70!black] {$\sigma_2$};
  \draw[future arrow=opochgreen!70] (s5) -- ++(2.0, 0)
    node[font=\tiny, above, pos=0.4, opochgreen!70!black] {$\sigma_2$};

  % Bend class 2 arrows up to UNIQUE
  \draw[future arrow=opochgreen!70] (s4.east) ++(2.0, 0) -- (unique.south west);
  \draw[future arrow=opochgreen!70] (s5.east) ++(2.0, 0) -- (unique.south west);

  % === Futures arrows from Class 3 states -> OMEGA ===
  \draw[future arrow=opochorange!70] (s6) -- ++(0.5, 0)
    node[font=\tiny, above, pos=0.5, opochorange!70!black] {$\sigma_1$};
  \draw[future arrow=opochorange!70] (s7) -- ++(0.5, 0)
    node[font=\tiny, above, pos=0.5, opochorange!70!black] {$\sigma_1$};
  \draw[future arrow=opochorange!70] (s8) -- ++(0.5, 0)
    node[font=\tiny, above, pos=0.5, opochorange!70!black] {$\sigma_1$};

  \draw[future arrow=opochorange!70] (s6.east) ++(0.5, 0) -- (omega.west);
  \draw[future arrow=opochorange!70] (s7.east) ++(0.5, 0) -- (omega.west);
  \draw[future arrow=opochorange!70] (s8.east) ++(0.5, 0) -- (omega.west);

  % === Annotation: equivalence meaning ===
  \node[font=\tiny, align=center, opochgray!70!black] at (4.0, -4.3)
    {States in the same class share identical futures:\\
     $s \MNcong s' \;\Leftrightarrow\; \forall \sigma:\;
     \mathrm{outcome}(s, \sigma) = \mathrm{outcome}(s', \sigma)$};

  % === Bidirectional equivalence arrows within classes ===
  \draw[{Stealth[length=1.5mm]}-{Stealth[length=1.5mm]}, densely dotted,
    opochblue!50, thin] (s1) -- (s2);
  \draw[{Stealth[length=1.5mm]}-{Stealth[length=1.5mm]}, densely dotted,
    opochblue!50, thin] (s2) -- (s3);

  \draw[{Stealth[length=1.5mm]}-{Stealth[length=1.5mm]}, densely dotted,
    opochgreen!50, thin] (s4) -- (s5);

  \draw[{Stealth[length=1.5mm]}-{Stealth[length=1.5mm]}, densely dotted,
    opochorange!50, thin] (s6) -- (s7);
  \draw[{Stealth[length=1.5mm]}-{Stealth[length=1.5mm]}, densely dotted,
    opochorange!50, thin] (s7) -- (s8);

\end{tikzpicture}
\caption{Myhill--Nerode congruence: states are partitioned into equivalence classes
  by futures equivalence $\MNcong$. States within the same class (dashed regions)
  produce identical terminal outcomes ($\UNIQUE$ or $\OMEGA$) under every test
  sequence~$\sigma$. The quotient automaton has exactly one state per class.}
\label{fig:myhill-nerode}
\end{figure}
